\caption{\textbf{Interval estimates for the reported quantities.} Seizure orders are not randomly assigned, so these are not tests of an intervention; they state how far each statistic would move under a different draw of the sample that produced it. Two units are reported. Resampling addresses treats the designated addresses as independent draws, which overstates the information available: addresses named in the same order share a date, an issuer and usually an operator. Resampling orders, and carrying their addresses with them, is the conservative reading, and only eight orders fall inside the observation window, so those intervals are wide. Where the two disagree, the order-level interval is the one to believe. Percentile intervals from 2,000 resamples, except the last row, which is a normal approximation over counterparties treated as independent and has no order-level analogue. The order-level column resamples the orders that contribute to each row: the eight orders inside the observation window for the timing rows, all twenty for freeze coverage, and the orders with at least one frozen address for the two frozen-address rows.}
\label{tab:uncertainty}
\small
\begin{tabular}{lrrrrr}
\toprule
Quantity & $n$ & Orders & Estimate & 95\% CI, addresses & 95\% CI, orders \\
\midrule
Median days from last transfer to signing & 178 & 8 & 32.5 & [9.5, 46] & [-77, 66] \\
Share dormant more than 30 days at signing & 178 & 8 & 0.511 & [0.438, 0.584] & [0.279, 0.764] \\
Share of designated addresses ever frozen & 490 & 19 & 0.808 & [0.773, 0.841] & [0.462, 0.973] \\
Frozen balance as a share of lifetime inflow & 396 & 16 & 0.0008 & [0.0004, 0.0015] & [0.0002, 0.0060] \\
Median days from last transfer to freeze & 396 & 16 & 117 & [82, 158] & [9, 227] \\
Share of counterparties transacting after the order & 20,295 & --- & 0.287 & [0.281, 0.293] & --- \\
\bottomrule
\end{tabular}
